\subsection{SUBRINGS; IDEALS; QUOTIENT RINGS; PRODUCTS OF RINGS}

If H is a subring of a topological ring A, then the topology induced on H by that of A is compatible with the ring structure of H. The topological ring structure thus defined on H is said to be induced by that of A.

PROPOSITION 5. Let H be a dense subring of a topological ring A, and let K be a subring (resp. left ideal, right ideal, two-sided ideal) of H. Then the closure K of K in A is a subring (resp. left ideal, right ideal, two-sided ideal) of A.

The proof is the same as for Proposition 8 of § 2, n.~3: if for example K is a left ideal in H, then the mapping \((z, x) \to zx\) is continuous on \(A \times A\) and maps \(H \times K\) into K; hence it maps \(A \times K = \overline{H} \times K\) into \(\overline{H}\) .

Let H be a two-sided ideal in a topological ring A. By the same argument as for quotient groups, we see that the quotient of the topology of A by the relation \(x - y \in H\) is compatible with the ring structure of A/H. In particular, if A is not Hausdorff, then the closure N of \(\{o\}\) in A is a closed two-sided ideal, by Proposition 5; the quotient ring, which is Hausdorff (§ 2, no. 5, Proposition 13) is called the Hausdorff ring associated with A.

Let \((\mathbf{A}_t)_{t\in \mathbf{I}}\) be a family of topological rings. On the set \(\mathbf{A} = \prod_{t\in \mathbf{I}}\mathbf{A}_t\) , the product of the topologies of the \(\mathbf{A}_t\) is compatible with the product of the ring structures of the \(\mathbf{A}_t\) (same proof as for product groups); the topological ring \(\mathbf{A}\) thus defined is called the product of the topological rings \(\mathbf{A}_t\) .

